\chapter{Клеточные автоматы, обратимость и представление Маделунга}
\label{app:background}

\section{Зачем это приложение}

Структурное определение клеточного автомата, дискретной фазы и координат $(dh,dt)$
дано в конце спуска (гл.~\ref{ch:descent}, \cref{sec:descent-ca,sec:coords-dh-dt}).
Здесь --- исторический и учебный контекст: классические имена, Тьюринг-полнота,
развёрнутый Маделунг и литература.
Читатель, идущий по основному тексту, может пропустить приложение.

\section{Исторические клеточные автоматы}

Идея восходит к \emph{Станиславу Уламу} и \emph{Джону фон Нейману} (1940-е):
сложное поведение из локальных правил на решётке.
\emph{«Игра жизни''} Конвея (1970) и \emph{Rule~110} (доказательство универсальности Cook, 2004)
показывают, что даже простые булевы правила Тьюринг-полны.
Наша модель~$M$ --- \emph{другой} элемент того же класса:
обратимый спинорный автомат с регистром $\mathbb{Z}_{512}[i]$, не подобранный для картинки.

\section{Тьюринг-полнота}
\label{sec:app-universality}

\emph{Тьюринг-полнота} означает способность имитировать машину Тьюринга при достаточном объёме.
Это не тезис «Вселенная --- ноутбук'', а методологический факт:
локальный автомат с конечным алфавитом \emph{достаточен} по классу динамик.
Абсолютные условия и геометрия FCC выделяют \emph{конкретный} закон~$g$ (гл.~\ref{ch:evolution}).

\begin{remark}[Три уровня]
\begin{enumerate}
\item \textbf{Класс CA} --- Тьюринг-полнота; локальность + дискретность.
\item \textbf{Закон~$g$} --- один обратимый элемент класса на FCC.
\item \textbf{Макроуровень~$T$} --- осреднение; Born \eqref{eq:born} на грубом приборе.
\end{enumerate}
\end{remark}

\section{Обратимость: Тоффоли и Фредкин}

На~$M$ требуется $g^{-1}(g(\Psi))=\Psi$ --- иначе унитарность \eqref{eq:norm} нарушается на шаге.
\emph{Эдвард Фредкин} и \emph{Томмазо Тоффоли} (1970--1980-е) развивали обратимые CA
как модель вычисления без энтропийного стирания бита.
Физическая мотивация в монографии та же, что в спуске:
микрошаг --- перестановка конечного множества состояний; необратимость измерения --- на~$T$.

\section{Представление Маделунга}

В 1927 году \emph{Эрвин Маделунг} переписал уравнение Шрёдингера в гидродинамической форме:
$\psi=\sqrt{\rho}\,e^{iS/\hbar}$, откуда неразрывность $\partial_t\rho+\nabla\cdot(\rho\mathbf{v})=0$
и квантовое давление $Q\propto\nabla^2\sqrt{\rho}/\sqrt{\rho}$.

Маделунг --- язык \emph{непрерывного}~$T$.
На~$M$ нет гладкого $\partial_t\rho+\nabla\cdot\mathbf{j}=0$ как фундаментального закона;
норма сохраняется дискретно, а гидродинамика появляется после осреднения
(гл.~\ref{ch:macro}, \ref{ch:matter}).
Дискретизация фазы в спуске (\cref{sec:descent-phase-disc}) --- не «доказательство Маделунгом'',
а замена $\varphi\in\mathbb{R}$ конечным $\varphi_{\mathrm{disc}}$ при том же смысле $|z|^2$ и фазы.

\begin{remark}[Частые заблуждения]
\begin{enumerate}
\item \textbf{«Маделунг доказывает дискретность.''} Нет --- он переписывает непрерывное уравнение.
\item \textbf{«КА = Игра жизни.''} Игра жизни --- частный необратимый пример.
\item \textbf{«На $M$ --- неразрывность Маделунга.''} Гладкость --- свойство~$T$ при $\lambda\gg\hl$.
\end{enumerate}
\end{remark}

\section{Сводная таблица}

\begin{table}[ht]
\centering
\caption{От классических понятий к главам монографии}
\label{tab:app-map}
\small
\begin{tabular}{@{}p{3.2cm}p{4.2cm}p{5.8cm}@{}}
\toprule
Понятие & Классика & У нас \\
\midrule
Клеточный автомат & Улама, von Neumann & Спуск \cref{sec:descent-ca}; закон~$g$ \\
Тьюринг-полнота & Конвей, Rule~110 & Класс CA; \S\ref{sec:app-universality} \\
Обратимость & Фредкин, Тоффоли & $g^{-1}$, гл.~\ref{ch:evolution} \\
Дискретная фаза & --- & Спуск \cref{sec:descent-phase-disc} \\
$(dh,dt)$ & --- & Спуск \cref{sec:coords-dh-dt} \\
Скорость света CA & $c_0=1$ cell/tick & $c_0=\hl/\htick$; см. \cref{sec:ca-terminology,sec:two-c} \\
Маделунг & 1927, $\rho$, $\mathbf{j}$ & Только~$T$ \\
\bottomrule
\end{tabular}
\end{table}

\section{Литература для первого чтения}

\begin{itemize}
\item E.~Fredkin, Digital mechanics (1990-е).
\item T.~Toffoli, N.~Margolus, \emph{Cellular Automata Machines} (MIT Press, 1987).
\item J.~von Neumann, \emph{Theory of Self-Reproducing Automata} (1966).
\item E.~Маделунг, Quantentheorie in hydrodynamischer Form (1927).
\item M.~Cook, Universality in elementary cellular automata (2004).
\item S.~Wolfram, \emph{A New Kind of Science} (2002) --- обзор; наша модель не следует его классификации правил.
\end{itemize}
